\documentclass[aspectratio=169,10pt]{beamer}
\usepackage[UTF8,fontset=none]{ctex}
\usepackage{booktabs}
\usepackage{array}
\usetheme{tutorial-red-shtu}

\title[Project Of SI100B EE Part]{Project Of SI100B EE Part}
\subtitle{RISCV miniCPU Design in Logisim}
\author{Chaofan Li}
\institute{ShanghaiTech University}
\date{}
\newcommand{\VSPsourcefooter}{~\emph{SI100B EE Part} \emph{Project Introduction} \emph{Dec.5 2025}}

\begin{document}
\VSPtitleframe

\begin{frame}{Contents}
  \tableofcontents
\end{frame}

\begin{frame}[fragile,shrink=12]{Background------What is CPU?}
\footnotesize
\begin{columns}[T,onlytextwidth]
  \begin{column}{.58\textwidth}
    \includegraphics[width=\linewidth,height=.54\textheight,keepaspectratio]{practice/Tutorial-SI100B-pj-intro/img/MC_CPU1.png}
  \end{column}
  \begin{column}{.38\textwidth}
    \includegraphics[width=\linewidth,height=.54\textheight,keepaspectratio]{practice/Tutorial-SI100B-pj-intro/img/MC_CPU2.png}
  \end{column}
\end{columns}
\end{frame}

\begin{frame}[fragile,shrink=12]{Background------Compile}
\footnotesize
\includegraphics[width=\linewidth,height=.54\textheight,keepaspectratio]{practice/Tutorial-SI100B-pj-intro/img/image.png}
\end{frame}

\begin{frame}[fragile,shrink=12]{Background------Assembly}
\footnotesize
\includegraphics[width=\linewidth,height=.54\textheight,keepaspectratio]{practice/Tutorial-SI100B-pj-intro/img/image-1.png}
\end{frame}

\begin{frame}[fragile,shrink=12]{Background}
\footnotesize
\begin{columns}[T,onlytextwidth]
  \begin{column}{.58\textwidth}
    \textbf{What is RISC-V?} 在计算机体系结构中，\textbf{指令集架构}（ISA）定义了处理器的基本指令规范，而\textbf{微架构}则负责具体实现这些指令的硬件电路设计。
    
    主流的指令集包括 x86、ARM、MIPS 等。与它们不同的是， RISC-V 作为一个全新的\textbf{开源指令集架构}正在迅速发展。与闭源的 x86 和 ARM 不同，RISC-V 采用开放标准，允许任何厂商自由使用和实现。
  \end{column}
  \begin{column}{.38\textwidth}
    \textbf{ISA vs. Microarchitecture}
    
    \begin{itemize}
    \item
      \textbf{ISA}: 软件能看到的指令和寄存器约定。
    \item
      \textbf{Microarchitecture}: 让这些指令真正运行起来的硬件组织。
    \item
      本 Project 的目标是用 Logisim 设计一个能执行 RISC-V 子集的 miniCPU。
    \end{itemize}
  \end{column}
\end{columns}
\end{frame}

\begin{frame}[fragile,shrink=12]{Overview of RISC-V Architecture}
\footnotesize
\includegraphics[width=\linewidth,height=.54\textheight,keepaspectratio]{practice/Tutorial-SI100B-pj-intro/img/RISCV-Arch.png}
\end{frame}

\begin{frame}[fragile,shrink=12]{Tasks}
\footnotesize
\textbf{Must Part} 1. 实现 miniCPU 的 Datapath，包括 RegFile、ALU、PC、IMem、DMem 等。（20\%）(在gradescope上有测试点) 2. 实现 miniCPU 的 Controller，包括指令译码、控制信号生成等，并与 Datapath 部分连接，实现指令执行。（20\%） 3. 使用 Python 编写一个简单的汇编器，将 RISC-V 汇编代码转换为机器码。（20\%） 4. 手动编写汇编程序，完成指定功能，加载到 miniCPU 的指令存储器中，观察 miniCPU 的运行情况。（15\%） 4. 实验报告。（25\%）
\end{frame}

\begin{frame}[fragile,shrink=12]{Tasks}
\footnotesize
\textbf{Must Part (Result Example)} \includegraphics[width=\linewidth,height=.54\textheight,keepaspectratio]{practice/Tutorial-SI100B-pj-intro/img/Result.png}
\end{frame}

\begin{frame}[fragile,shrink=12]{Bonus}
\footnotesize
\begin{enumerate}
\def\labelenumi{\arabic{enumi}.}
\item
  为单周期 miniCPU 添加流水线。 \textgreater{} 添加流水线之后，你可能会遇到被称之为``冒险''（Hazard）的问题，你需要设计合适的解决方案来处理这些冒险问题。
\item
  为 miniCPU 增加更多指令的支持。
\end{enumerate}

\begin{vspquote}{Bonus Challenge}
我觉得自己特别厉害！那么可以尝试以下 Bonus 3. 为 miniCPU 添加向量运算单元，并支持一些向量指令的支持。 4. 为 miniCPU 支持标量-向量双发射。 5. 使用 Spinal HDL 完成 miniCPU 的设计与实现！（\textbf{Recommended}） - 如果你决定使用 Spinal HDL 来实现 miniCPU，你可以不实现 Logisim 版本。请务必提前联系助教！完成相同的任务，使用 Spinal HDL 实现的项目将获得额外的加分。
\end{vspquote}
\end{frame}

\section{Grading Policy}

\begin{frame}[fragile,shrink=12]{About AI?}
\footnotesize
\begin{itemize}
\item
  在我们的 Project 中，我们\textbf{允许任何 AI 工具的使用}。但是，请\textbf{务必确保你理解你的 AI 做了什么}。如果在最终的 Check 中你的回答是``AI 这么写的，我也不太清楚''，那么我们可能会认为你并没有真正理解你的实现，从而影响你的评分。
\item
  请同学们在合作中遵守基本的学术诚信原则。\textbf{组间的抄袭}是严格禁止的。
\item
  允许参考公开的资源实现，但请按照自己的理解进行调整和修改。如果你参考了公开的资源（例如 GitHub 上的开源项目），请务必在报告中注明引用来源。非复制粘贴的参考不会影响你的评分。 
\includegraphics[width=\linewidth,height=.54\textheight,keepaspectratio]{practice/Tutorial-SI100B-pj-intro/img/99e4b731e1a0ee359087f8ed1151dea4.jpg}
\item
  即使存在 Bug，如果你能准确描述 Bug 的原因，并展示你为解决 Bug 所做的努力，我们也会给予一定的分数。实现固然重要，但我们更看重你在实现过程中所展现的理解和思考！
\end{itemize}
\end{frame}

\begin{frame}[fragile,shrink=12]{Grading Policy for Bonus Tasks}
\footnotesize
\textbf{Bonus} - Bonus 部分的分数将在基础分数的基础上进行加分。 - 这部分不会有 OJ 测试，但我们会在面对面 Check 的时候，检查你们的实现。 - 你们需要自己设计合适的测试用例，以证明你们的实现是正确的。否则，我们可能会认为你们的实现是不完整的，从而你们的加分也会受到影响。 - 和 Must 一样，即使你们的实现不完整，如果在报告中展现了你们的尝试和思考，并做出有意义的尝试和分析，我们也会给予一定的加分。
\end{frame}

\begin{frame}[fragile,shrink=12]{Who should choose our project?}
\footnotesize
\begin{itemize}
\item
  对``计算机到底是怎么运行程序的''感兴趣的同学

  \begin{itemize}
  \item
    不要求你来自信院，也不要求你已经学过体系结构或数字电路
  \item
    只要你愿意把一个系统拆开来看，理解从指令、寄存器到简单 CPU 的基本工作过程
  \end{itemize}
\item
  来自生命、物质、创艺与艺术等专业，但希望补一块计算机底层视角的同学

  \begin{itemize}
  \item
    生命/生技方向：如果你未来会接触计算、生物信息、自动化实验或数据分析，理解计算机如何执行程序会很有帮助
  \item
    物理/化学/材料方向：如果你关心仿真、仪器控制、高性能计算或硬件相关工具，这个 Project 可以帮你建立硬件和软件之间的连接
  \item
    创艺方向：如果你对交互装置、创意编程、生成式艺术或软硬件结合感兴趣，这里会提供一个更接近机器本身的入口
  \end{itemize}
\item
  希望挑战一点``看起来不像本专业''的内容，但不想直接跳进高 workload 专业课的同学

  \begin{itemize}
  \item
    我们会更看重你是否理解自己的实现，而不是你一开始是否有相关背景
  \item
    即使最后没有做得非常完整，只要能讲清楚尝试、问题和思考，也会有相应的评价空间
  \end{itemize}
\end{itemize}
\end{frame}

\VSPendframe{SI100B EE Part \textcolor{white!60}{/} Project Introduction \textcolor{white!60}{/} Q\&A}
\end{document}
